% Sections 1-2 (Introduction, Background), written 21 Oct 2026.
% NOT yet in main.tex: to be inserted together with Section 5.

\section{Introduction}\label{sec:intro}

Every runner who follows a training plan is, without knowing it, controlling a
dynamic system. The input is the training done every day; the output is the
performance on race day. In between there are two quantities that cannot be
measured directly: the fitness built up over weeks and months, and the fatigue
accumulated in the last few days. Hard training raises both. Fitness lasts
long, fatigue goes away quickly, and the performance on a given day depends on
the balance between the two. This is the reason why runners reduce their
training in the last days before a race (the taper): fatigue disappears first,
most of the fitness remains.

Coaches manage this balance with experience. The idea of this project is to
manage it as an engineer: write a model of the athlete, identify its
parameters from data, estimate its hidden states and use it to plan the
training. The athlete is the author, and the data are the ones recorded by his
own sports watch during almost one year of running. The case study is the
preparation of a half marathon on 22~November~2026: the system is used while
the preparation is still running, and the race is its final test.

The project follows the classical chain of a control problem:
\begin{enumerate}
  \item \textbf{Data and input.} A pipeline converts the raw files of the
  watch into a daily training load, computed from heart rate second by second
  (Sections~\ref{sec:data} and~\ref{sec:load}).
  \item \textbf{Output.} Since races are rare, performance is measured on
  ordinary training runs through an aerobic-efficiency index, whose noise is
  quantified and whose dependence on weather and place is checked
  (Section~\ref{sec:perf}).
  \item \textbf{Identification.} The parameters of the fitness--fatigue model
  are estimated from these data, taking advantage of three intentional rest
  periods that act as natural step experiments (Section~\ref{sec:id}).
  \item \textbf{State estimation.} A Kalman filter estimates fitness and
  fatigue day by day as new runs are recorded (Section~\ref{sec:kf}).
  \item \textbf{Planning.} Model predictive control computes the optimal
  training load for the last weeks before the race, and the result is compared
  with the plan actually followed (Section~\ref{sec:mpc}).
  \item \textbf{Test.} The predicted race time is committed to the code
  repository before the race, with a date that cannot be changed, and then
  compared with the real result (Section~\ref{sec:race}).
\end{enumerate}

The last point is what makes the project a real test and not only a fit of
past data: a model adapted after seeing the result can always be made to work,
a prediction fixed in advance cannot. For the same reason, until the race the
model is used only as an advisor, and the training plan is not changed on the
basis of its output.

% ---------------------------------------------------------------------
\section{Background and related work}\label{sec:background}

\subsection{The fitness--fatigue model}
The first systems model of training was proposed by \citet{banister1975} and
described in detail by \citet{calvert1976}. The athlete is treated as a linear
system whose impulse response is the difference of two exponentials: a single
training session produces a positive effect (fitness), large and slow, and a
negative effect (fatigue), smaller in the long run but faster and initially
stronger. Discretised on a daily grid, with training load $w[k]$, the model
reads:
\begin{align}
x_1[k+1] &= e^{-1/\tau_1}\, x_1[k] + w[k] && \text{(fitness)}\label{eq:x1}\\
x_2[k+1] &= e^{-1/\tau_2}\, x_2[k] + w[k] && \text{(fatigue)}\label{eq:x2}\\
p[k]     &= p_0 + k_1 x_1[k] - k_2 x_2[k]  && \text{(performance)}\label{eq:p}
\end{align}
where $\tau_1 > \tau_2$ are the time constants in days, $k_1$ and $k_2$ the
gains and $p_0$ the baseline performance. With the typical values of
$\tau_1 \approx 40$ days and $\tau_2 \approx 7$ days, a single hard session
first makes the athlete worse and only after some days better: this is the
mechanism behind both recovery and taper.

From the engineering point of view, the model is a linear time-invariant
system with one input, two states and one output, and can be written in the
standard state-space form
\begin{equation}
\mathbf{x}[k+1] = \mathbf{A}\,\mathbf{x}[k] + \mathbf{B}\,w[k], \qquad
p[k] = \mathbf{C}\,\mathbf{x}[k] + p_0,
\label{eq:ss}
\end{equation}
with $\mathbf{A} = \mathrm{diag}(e^{-1/\tau_1}, e^{-1/\tau_2})$,
$\mathbf{B} = [1\;\;1]^{\mathsf T}$ and $\mathbf{C} = [k_1\;\;-k_2]$.
This form is the reason why the model fits so naturally with the tools used in
this project: identification, Kalman filtering and model predictive control
all start from a state-space model.

\subsection{Applications and limitations}
\citet{morton1990} applied the model to running performance, and
\citet{fitzclarke1991} used it to find how training sessions should be placed
before a competition, which is the first example of planning based on the
model. Later work focused on its limits. \citet{busso1997} estimated the
parameters recursively, letting them change over time, and \citet{busso2003}
showed that the fatigue response grows with the recent training load, so that
the same session costs more when the athlete is already tired.
\citet{hellard2006} analysed the model on elite athletes and found that its
parameters are estimated with large uncertainty and are strongly correlated
with each other, so that their values must be interpreted with care. Nonlinear
alternatives have also been proposed, such as the PerPot model of
\citet{perl2001} and the model of \citet{turner2017}, which was also used to
optimise the training plan. A clear introduction to the whole topic, with
worked examples, is given by \citet{clarke2013}.

\citet{kolossa2017} combined the fitness--fatigue model with a Kalman filter,
using new performance measurements to correct the state of the model as they
become available. This is the closest work to the state-estimation part of the
present project.

\subsection{Measuring the training load}
The input of the model must be a single number per day. \citet{banister1991}
proposed the training impulse (TRIMP), which weights the duration of the
exercise by the heart-rate reserve with an exponential factor, so that minutes
at high intensity count more than minutes at low intensity. TRIMP is a
measure of \emph{internal} load, i.e.\ of the response of the body, while
distance and pace are measures of \emph{external} load \citep{impellizzeri2019}.
For a model of the response of the athlete, the internal load is the natural
input.

\subsection{Position of this work}
Most studies on the fitness--fatigue model use data collected in laboratory
tests or races, on groups of athletes, and fit the model a posteriori. This
project differs in three respects: it uses only the data of a consumer
sports watch recorded during normal training; it measures performance on
ordinary training runs, so that measurements are frequent; and it uses the
model in real time during a preparation, with a prediction fixed before the
race. The control tools themselves are standard and are taken from the
textbooks of \citet{ljung1999} for identification, \citet{simon2006} for
state estimation and \citet{rawlings2017} for model predictive control.
